% from ega3-1-fr.tex
\oldpage[0\textsubscript{III}]{72}
Ces trois propriétés résultent aussitôt des définitions de
(\hyperref[0.11.2.2-fr]{11.2.2}), en tenant compte de ce que
$\mathrm F^{p-r+1}(X_k^\bullet)=X_k^\bullet$ pour $p-r<p_0$.

\begin{env}[13.5.4]
\label{0.13.5.4-fr}
Supposons vérifiées les hypothèses de
(\hyperref[0.13.5.3-fr]{13.5.3}). Alors, pour $p,q,r$ fixés ($r$ fini), les
systèmes projectifs
\[
(\mathrm Z_r^{pq}(X_k^\bullet))_{k\in\mathbb Z},\quad
(\mathrm B_r^{pq}(X_k^\bullet))_{k\in\mathbb Z},\quad
(\mathrm E_r^{pq}(X_k^\bullet))_{k\in\mathbb Z}
\]
sont \emph{essentiellement constants}~; on désignera par
$\mathrm Z_r^{pq}(\mathbf X^\bullet)$,
$\mathrm B_r^{pq}(\mathbf X^\bullet)$ et
$\mathrm E_r^{pq}(\mathbf X^\bullet)=
\mathrm Z_r^{pq}(\mathbf X^\bullet)/
\mathrm B_r^{pq}(\mathbf X^\bullet)$ leurs limites projectives
respectives. Les $\mathrm Z_r^{pq}(\mathbf X^\bullet)$ et
$\mathrm B_r^{pq}(\mathbf X^\bullet)$ s'identifient canoniquement à des
sous-objets de $\mathrm E_1^{pq}(\mathbf X^\bullet)$. La définition des
$d_r^{pq}$ (M, XV, 1) montre que ces morphismes (relatifs aux
$X_k^\bullet$) sont aussi essentiellement constants, et par suite
définissent des morphismes
\begin{equation}
\label{0.13.5.4.1-fr}
\tag{13.5.4.1}
d_r^{pq}:\mathrm E_r^{pq}(\mathbf X^\bullet)
\longrightarrow\mathrm E_r^{p+r,q-r+1}(\mathbf X^\bullet)
\end{equation}
tels que $d_r^{p+r,q-r+1}\circ d_r^{pq}=0$~; en outre, on a des
isomorphismes canoniques de $\operatorname{Ker}(d_r^{pq})$ sur
$\mathrm Z_{r+1}^{pq}(\mathbf X^\bullet)/
\mathrm B_r^{pq}(\mathbf X^\bullet)$ et de
$\operatorname{Im}(d_r^{pq})$ sur
$\mathrm B_{r+1}^{p+r,q-r+1}(\mathbf X^\bullet)/
\mathrm B_r^{p+r,q-r+1}(\mathbf X^\bullet)$.
\end{env}

\begin{lemma}[13.5.5]
\label{0.13.5.5-fr}
\emph{Sous les hypothèses de
(\hyperref[0.13.5.3-fr]{13.5.3}), on a, pour $s\geq r>p-p_0$, un
monomorphisme canonique}
\begin{equation}
\label{0.13.5.5.1-fr}
\tag{13.5.5.1}
i:\mathrm E_s^{pq}(\mathbf X^\bullet)
\longrightarrow\mathrm E_r^{pq}(\mathbf X^\bullet)
\end{equation}
\emph{et un isomorphisme canonique}
\begin{equation}
\label{0.13.5.5.2-fr}
\tag{13.5.5.2}
j_r:\mathrm E_r^{pq}(\mathbf X^\bullet)
\xrightarrow{\sim}\mathrm E_\infty^{pq}(X_{p+r-1}^\bullet)
\end{equation}
\emph{tels que le diagramme}
\begin{equation}
\label{0.13.5.5.3-fr}
\tag{13.5.5.3}
\xymatrix{
\mathrm E_s^{pq}(\mathbf X^\bullet)\ar[r]^{j_s}\ar[d]_i
&\mathrm E_\infty^{pq}(X_{p+s-1}^\bullet)\ar[d]\\
\mathrm E_r^{pq}(\mathbf X^\bullet)\ar[r]^{j_r}
&\mathrm E_\infty^{pq}(X_{p+r-1}^\bullet)
}
\end{equation}
\emph{soit commutatif} (la flèche verticale de droite provenant du
morphisme $X_{p+s-1}^\bullet\to X_{p+r-1}^\bullet$).
\end{lemma}
\begin{proof}
L'existence de $i$ provient de ce que
$\mathrm B_r^{pq}(X_k^\bullet)=\mathrm B_\infty^{pq}(X_k^\bullet)$ pour
$r>p-p_0$ (\hyperref[0.13.5.3-fr]{13.5.3}, a)~; on a
$\mathrm Z_r^{pq}(X_k^\bullet)=\mathrm Z_\infty^{pq}(X_k^\bullet)$ pour
$k+1\leq p+r$ (\hyperref[0.13.5.3-fr]{13.5.3}, b), d'où en particulier
$\mathrm Z_r^{pq}(X_{p+r-1}^\bullet)=
\mathrm Z_\infty^{pq}(X_{p+r-1}^\bullet)$ et d'autre part
$\mathrm Z_r^{pq}(X_{p+r-1}^\bullet)$ et
$\mathrm B_r^{pq}(X_{p+r-1}^\bullet)$ s'identifient canoniquement à
$\mathrm Z_r^{pq}(\mathbf X^\bullet)$ et
$\mathrm B_r^{pq}(\mathbf X^\bullet)$ en vertu de
(\hyperref[0.13.5.3-fr]{13.5.3}, c)), d'où l'existence de $j_r$ et la
commutativité de
(\hyperref[0.13.5.5.1-fr]{13.5.5.1}).
\end{proof}

\begin{corollary}[13.5.6]
\label{0.13.5.6-fr}
\emph{Sous les hypothèses de
(\hyperref[0.13.5.3-fr]{13.5.3}), si l'une des limites projectives
$\varprojlim_r\mathrm E_r^{pq}(\mathbf X^\bullet)$,
$\varprojlim_k\mathrm E_\infty^{pq}(X_k^\bullet)$ existe, il en est de
même de l'autre, et on a un isomorphisme canonique}
\begin{equation}
\label{0.13.5.6.1-fr}
\tag{13.5.6.1}
j_\infty:\varprojlim_r\mathrm E_r^{pq}(\mathbf X^\bullet)
\xrightarrow{\sim}
\varprojlim_k\mathrm E_\infty^{pq}(X_k^\bullet).
\end{equation}
\emph{En outre, pour que le système projectif
$(\mathrm E_r^{pq}(\mathbf X^\bullet))_{r\in\mathbb Z}$ soit
essentiellement constant
(\hyperref[0.13.4.2-fr]{13.4.2}), il faut et il suffit que le système
projectif $(\mathrm E_\infty^{pq}(X_k^\bullet))_{k\in\mathbb Z}$ le soit.}
\end{corollary}

\begin{env}[13.5.7]
\label{0.13.5.7-fr}
On note $\mathrm B_\infty^{pq}(\mathbf X^\bullet)$ et
$\mathrm Z_\infty^{pq}(\mathbf X^\bullet)$ les sous-objets de
$\mathrm E_1^{pq}(\mathbf X^\bullet)$ égaux respectivement à
$\mathrm B_r^{pq}(\mathbf X^\bullet)$ pour $r>p-p_0$ et à
$\inf_r\mathrm Z_r^{pq}(\mathbf X^\bullet)$ (quand ce dernier existe), de
sorte que $\varprojlim_r\mathrm E_r^{pq}(\mathbf X^\bullet)$
